 
\documentclass[a4paper,12pt]{article}
\usepackage[utf8]{inputenc}
\usepackage[margin=2.5cm]{geometry}
%\usepackage{geometry}[2cm,2cm,2cm,2cm]
\usepackage[utf8]{inputenc}
\usepackage[nosumlimits]{amsmath}
\usepackage{amssymb,stmaryrd,amsthm,MnSymbol}
\usepackage{amsfonts,amssymb,amsopn}
\usepackage[colorlinks=true,
            linkcolor=red,
            filecolor=blue,
            citecolor=blue,
            urlcolor=blue,]{hyperref}
\usepackage{color}
\usepackage[dvipsnames]{xcolor}
\usepackage{comment}
\usepackage{graphicx}
\usepackage{natbib}
\bibliographystyle{elsarticle-harv}
\usepackage{multirow}

\usepackage{soul}
\setstcolor{red}
\newcommand{\werner}[1]{{\color{magenta}WB says: #1}}

\newcommand{\todo}[1]{\vspace{5 mm}\par \noindent
	\framebox{\begin{minipage}[c]{0.95 \textwidth}
			\tt #1 \end{minipage}}\vspace{5 mm}\par}


%opening
\title{Response to reviewers}
\author{Werner Bauer, Colin J. Cotter}

\begin{document}

\maketitle

\paragraph{TODO}

\begin{enumerate}
\item 4/8 core runs for IMEX/GL1/GL2/GL3, and update table
\item Do code release and Zenodo
\item Redo plots in figure 3 without background
\end{enumerate}

We would like to thank the Reviewers for the constructive comments
which helped us to improve the manuscript.

In preparing our response, we noticed a couple of things that have led
to changes in our results. First, we realised that despite describing
the vector invariant form in the exposition, our results actually used
the advective form (due to us incorrectly selecting a switch in the
code).  The vector invariant form is more suited to the sphere, and we
anyway intended to use it, so we have rerun all our results using
vector invariant form, leading to slightly different errors. Second,
we realised that the IMEX and IRK results were not converging to each
other as $\Delta t\to 0$. This turned out to be a problem with a
tolerance in the IMEX mass matrix solve being too loose, so it was
effectively solving a different ODE. We have changed this, and it now
leads to slightly slower IMEX results (the solve is not dominated by
the mass solve but it does add a little time). We apologise to the Editor
and Reviewer for these inaccuracies in the first manuscript.

Additionally, we have improved the treatment of the Eisenstat-Walker
scheme, by adding an automated absolute tolerance to the Krylov method
which avoids over-solution of the linear system in the final Newton
iteration beyond the tolerance of the Newton solver. We have added 
\werner{sentence not complete here!}

The questions and comments by the Reviewer are written in Roman font,
followed by our responses in italic fonts. In addition, we marked in
the revised manuscript those parts in blue (Reviewer 1) and magenta
(Reviewer 2), which have been rewritten in order to highlight the
changes. Our own changes are marked in green.

\section{Reviewer 1}

%% The work present an implicit Runge Kutta method for solving the
%% rotating SW equations, focusing on how to solve the nonlinear system
%% with a quasi Newton method and on the preconditioners used to
%% preprocess the solution of the monolithic system of equations. In
%% particular, it uses an Additive Schwarz method to split the domain in
%% components to get an initial guess, multigrid nested meshes and Krylov
%% subspace preconditioner.

%% The work compares this algorithm with an IMEX RK on a test of a
%% nonlinear Rossby wave solution.

%% \subsection{Comment}
%% The work is well written and mathematically correct on what
%% presents. Nevertheless, the proposed algorithm is composed of many
%% ingredients that overlaps and get compared in the description section
%% so that it is not easy to have a full picture of its structure by
%% reading that section. Also, the description of the finite element
%% method, of the algorithm and of the test under consideration lack many
%% details that could help the reader and are left to the references. I
%% think adding some details would help understanding the work and
%% wouldn't lengthen it too much. Finally, the comparison with another
%% algorithm is not too fair as they scale very different in terms of
%% stability and accuracy. It is really hard to compare them and to
%% decide if the novel algorithm is worth the effort of a more
%% complicated fully implicit structure. Other tests could help validate
%% this aspect. I list below some remarks that should be addressed before
%% publication.

%% \subsection{Remarks}
\begin{enumerate}
\item page 3: Define the acronym BDM \\
  \emph{Thank you, this has been done.}

\item Page 3 eq 5: shouldn't it be g(D-b) according to eq
  1?\\ \emph{In fact, Eq. 1 was wrong, we have used the other sign
  convention (positive $b$ means raised topography) in the code and
  throughout the paper otherwise, so this has now been fixed.}
\item Page 3: the notation $[[ \cdot ]]$ for the average is uncommon
  and reminds more of jump rather than an average \emph{In fact, this
  does result in jumps when $n$ is present, because it changes sign on
  each side of the facet. There are various ways to do this in the literature,
  and we have chosen one that is not uncommon.
  We have added a footnote to clarify this.}
  
\item Page 4, last sentence should be revisited, it is complicated to read.
\emph{We have done this.}
  
\item Section 2.2 is somewhat chaotic, there are many different
  ingredients listed in the section, but it is difficult to follow the
  precise algorithm the authors use to solve their IRK problem. It
  would be useful to put an algorithmic box for the solution of the
  linear solver that explains the order of the algorithms, and it
  should be more clear what has just been compared and what has been
  used in the end.\\
  \emph{This is a very good suggestion, we have done this.}

\item Page 6: the sentence "For this type of problem, it is also
  typical to scale $\Delta t$ to keep the advective Courant number
  ($U(\Delta t/\Delta x) $ where U is a typical velocity scale and
  $\Delta x$ is a typical cell diameter) constant as the mesh is
  refined, so we can hope for mesh independent Krylov iteration counts
  under this refinement scaling without multigrid. However, we found
  that using the multigrid produced faster wallclock times." was to me
  clear only at page 8 when the authors state that IMEX are restricted
  by the Courant number on the explicit part and that "means that it
  is limited to smaller timesteps. It is this property that we shall
  contrast with the IRKs." Make clearer from the beginning that you
  are using more iterations and larger timesteps.\\
  \emph{We agree that this was not very clear. We have now clarified
  that this means that using multigrid is faster than just using
  a patch smoother on the finest grid.}
\item There are many tolerances in the proposed algorithm that make the comparison not so fair, it would be interesting to see if it is stable under perturbation of these tolerances.\\
{\itshape
  We would like to emphasise that the comparison with IMEX is
  not to try to compete with it. We include IMEX results as a point of
  reference so that the reader can decide if IRKs are viable using our
  solver approach. The idea is that IMEX schemes are commonly used in
  the community and so the reader might have a feel for how long IMEX
  schemes take for their preferred discretisation. We have clarified this
  in the text.

  Regarding the tolerances, there is really only one tolerance in the
  IRK methodology, which is the relative tolerance for the Newton
  iterations.
  We have included a discussion of the effects of changing this tolerance in
  the text.
  
  As we mentioned above, the tolerance in the mass matrix solve for the
  IMEX scheme is important to keep tight, since otherwise we do not converge
  to the right solution as $\Delta t \to 0$. However, this mass solve
  does not dominate the IMEX cost.}
\item Comparing algorithms to only one test case could be
  misleading. In fact the errors of the IMEX and of the IRK (low
  order) for this test are so different that is really impossible to
  compare them.\\
  \emph{The goal here is not to say that IRKs are better or worse than IMEX
  schemes in general, but just to give the reader a point of reference
  as written above. In any case, the timings have changed quite a lot after
  our updates.}
\item I also suggest to find another algorithm to compare the proposed IRK method as the IMEX is so much more accurate and already unstable with relatively small timesteps that the curves are so far apart. Either another test that makes the two more comparable or another state-of-the-art algorithm (a DIRK for example)?\\
  \emph{We have now included a 3rd order 3 stage DIRK, for comparison.
  However, note that this DIRK also required a scalable solver strategy,
  we used the same one as we used for the IRKs.
  }
\end{enumerate}




\section{Reviewer 2}

%% The authors present some preliminary results on the application of
%% fully-implicit Runge-Kutta (Gauss and Radau-IIA) methods to a finite
%% element discretization of the rotating shallow water equations on the
%% sphere. They combine some existing tools and approaches, and show some
%% numerical results including a comparison with a second-order IMEX
%% method. The results can be interesting for the community, but I would
%% like to see several improvements before recommending publication.



\subsection{Remarks}
\begin{enumerate}
\item  Abstract, "Second, particular choices of implicit timestepping methods can extend energy conservation properties of spatial discretisations to the fully discrete method." This is not the case for the methods used in this paper.

  \emph{We have removed this from the abstract, as this kind of discussion is
  anyway not really suitable for an abstract.}
  
\item  Abstract, "Third, these methods avoid issues related to splitting errors that can occur in some situations, and avoid the complexities of splitting methods." Instead, fully implicit methods require solving nonlinear systems with Jacobian changing at each iteration. Thus, the complexity increases significantly.

  \emph{Similarly, we have removed this from the abstract.}
  
\item  It would be helpful for readers not familiar with the notation to explain the significance of the semicolon in the forms $a, c$ defined in (3) and (4).

\emph{Thank you, this has been done.}
  
\item  References to concrete results in books have to be more specific, e.g., by mentioning the section/theorem/page number (e.g., on page 4).

\emph{We have added some details about the relevant chapters in this book citation.}
  
\item Page 4: The time derivatives $k_u$ and $k_D$ are used as primary
  variables in the Newton iteration. However, in their well-known
  textbook, Hairer and Wanner recommend using updates of the stage
  values $u$ and $D$ instead (Section IV.8). Please comment on the
  choice, explain, and discuss the advantages/disadvantages.\\
  {\itshape It does not seem to be the case that Hairer and Wanner
  recommend using the stage values, it is just that they only discuss
  this formulation in the implementation section. In that textbook,
  they seem to presume that you must avoid full Newton and do
  something cheaper in practice, but that's not really the case with
  modern computational implementations such as PETSc.

  We made a preliminary investigation of the stage values formulation,
  since this option is easily available in Irksome. However, we
  obtained linear solver divergences on the second Newton iteration
  for the largest timestep considered (7200s), which suggests that a
  better initial guess is needed at large timesteps (we used the value
  of the solution at the start of the timestep for all stages). We did
  not encounter this problem for stage derivatives, even for larger values
  not presented in the manuscript. For smaller values there is no difference
  in performance or error for stage values, so this does not seem to be
  an important aspect.}
\item  Continuing the previous comment: Hairer and Wanner also recommend transforming the linear systems to solve multiple smaller systems instead of one large system (in the same section). Please comment on this as well.

  \emph{Hairer and Wanner assume a direct solver is being used for the
  entire system, whilst we are using an iterative solver. Solving one
  large system versus several small systems will only make a
  difference if the preconditioner behaviour differs. We do have
  direct solvers for the patch systems, but at this point it becomes
  dangerous to think only of operation counts since the bottleneck can
  frequently be the time required to stream RHS data to the patch
  solves.  Hence it is probably safer to decline to comment on this
  without doing a proper performance model and benchmarking of the
  patch solves.  }
  
\item  Page 4: Please explain in detail what "sufficiently small" means.
  \emph{We have described the relative tolerance test used.}
  
\item Please provide numerical results concerning the long-term
  stability of the methods for more demanding test cases. For example,
  run the Rossby-Haurwitz test case for 28 days instead of one day,
  and consider also other test cases involving instabilities.

\emph{These have now been included. For instabilities we ran the
Williamson test case 5 to 50 days, after which it has developed a lot of
vorticity structure.}
  
\item  Aims and scope of the journal: "Reproducibility, that is the ability to reproduce results obtained by others, is a core principle of the scientific method. As the impact of and knowledge discovery enabled by computational science and engineering continues to increase, it is imperative that reproducibility becomes a natural part of these activities. The journal strongly encourages authors to make available all software or data that would allow published results to be reproduced and that every effort is made to include sufficient information in manuscripts to enable this. This should not only include information used for setup but also details on post-processing to recover published results."

\emph{We have now addressed this by referencing our code and providing a DOI
  to the versions of Firedrake/PETSC/etc used.}
\end{enumerate}


\end{document}
